\documentclass[xcolor=dvipsnames, aspectratio=169, 10pt]{beamer}

% ── Beamer Theme ──
\usetheme{Madrid}
\usecolortheme{whale}
\setbeamertemplate{navigation symbols}{}
\setbeamertemplate{footline}[frame number]

% ── Encoding & Fonts ──
\usepackage[utf8]{inputenc}
\usepackage[T1]{fontenc}
\usepackage[french]{babel}
\usepackage{inconsolata}
\usepackage{booktabs}
\usepackage{xcolor}
\usepackage{tabularx}

% ── Couleurs ──
\definecolor{primaryblue}{RGB}{0, 82, 164}
\definecolor{accentorange}{RGB}{230, 126, 34}

% ── Métadonnées ──
\title{\textbf{LockBits — Solution EDR Open-Source\\\& Portail Client}}
\subtitle{Dossier d'Architecture — Présentation Orale}
\author{Cléry A.-Ferradou \and Swann Kechar \and Soltane Afellah \and Quentin Labourdette}
\institute{ESGI — Projet Annuel 2\ieme{} année — 2025-2026}
\date{10 Juillet 2026}

% ── Helpers ──
\newcommand{\bitem}{\item[$\bullet$]}

\begin{document}

% ========================================================================
% SLIDE 1 — TITLE
% ========================================================================
\begin{frame}
  \titlepage
  \vfill
  \centering
  \small\textbf{Équipe}\\[2pt]
  Cléry A.-Ferradou \textit{(Chef de Projet / DevOps / Hébergeur)} —
  Swann Kechar \textit{(EDR)} —
  Soltane Afellah \textit{(Site Client)} —
  Quentin Labourdette \textit{(Chatbot)}
\end{frame}
\note{
  20 minutes, 4 parties.
  Cléry : contexte, architecture, CI/CD, hébergement.
  Swann : EDR serveur + agent.
  Soltane : site client.
  Quentin : chatbot + retour d'expérience général.
}

% ========================================================================
% SECTION CLÉRY — Équipe, Contexte, Architecture, CI/CD, Infrastructure
% ========================================================================
\section{Cléry — Organisation \& Infrastructure}

% --- Slide 2 : Organisation de l'équipe ---
\begin{frame}{Organisation de l'Équipe}
  \centering
  \begin{tabularx}{\textwidth}{@{}l X@{}}
    \toprule
    \textbf{Membre} & \textbf{Rôle \& Responsabilités} \\
    \midrule
    Cléry A.-Ferradou &
    \textbf{Chef de Projet / DevOps / Hébergeur} —
    Coordination, CI/CD, Docker, Serveur Proxmox, Cloudflare, \\
    Swann Kechar &
    \textbf{EDR} — Serveur FastAPI, Agent Python, VirusTotal, GLPI, \\
    Soltane Afellah &
    \textbf{Site Client} — PHP/Apache, MySQL, GLPI SSO, Tailwind CSS, \\
    Quentin Labourdette &
    \textbf{Chatbot} — Claude API, Sécurité, Support client, \\
    \bottomrule
  \end{tabularx}
  \vspace{0.3cm}
  \textbf{Répartition des dépôts Git}
  \begin{itemize}
    \bitem \texttt{main} : Cléry — Orchestrateur, CI/CD, documentation LaTeX
    \bitem \texttt{edr} : Swann — Serveur + Agent EDR
    \bitem \texttt{site\_lockbits} : Soltane + Quentin — Site + Chatbot
  \end{itemize}
\end{frame}
\note{
  L'équipe est organisée par composant : chacun son domaine.
  Cléry gère l'infrastructure commune (CI/CD, serveur, déploiement).
  Swann : le cœur technique (EDR).
  Soltane : l'interface client.
  Quentin : le support utilisateur (chatbot).
  Communication via GitHub Issues, PRs, Discord.
}

% --- Slide 3 : Contexte & Objectifs ---
\begin{frame}{Contexte \& Objectifs de l'Infrastructure}
  \textbf{Pourquoi LockBits ?}
  \begin{itemize}
    \bitem Entreprise fictive de \textbf{cybersécurité}
    \bitem Objectif : solution complète de détection de menaces + portail client + assistant
    \bitem \textbf{Contrainte} : 4 étudiants, 8 mois, matériel personnel, zéro budget
  \end{itemize}
  \vspace{0.2cm}
  \textbf{Ce qu'on a livré}
  \begin{itemize}
    \bitem \textbf{EDR} — Agent de détection + Serveur d'analyse + Alerting
    \bitem \textbf{Portail client} — Dashboard + Tickets GLPI + Téléchargement agent
    \bitem \textbf{Chatbot} — Assistant virtuel basé sur Claude API
    \bitem \textbf{Infrastructure} — CI/CD automatisé, Docker, Monitoring, Hébergement 24/7
  \end{itemize}
\end{frame}
\note{
  LockBits est une société fictive créée pour le projet annuel.
  L'objectif pédagogique : couvrir tout le cycle de vie d'un projet IT — du code au déploiement, en passant par l'infrastructure et la documentation.
  On bosse avec du matériel perso : serveur Proxmox chez Cléry.
}

% --- Slide 4 : Choix Techniques ---
\begin{frame}{Choix Techniques \& Stack}
  \centering
  \begin{tabularx}{\textwidth}{@{}l X@{}}
    \toprule
    \textbf{Composant} & \textbf{Technologie retenue} \\
    \midrule
    Serveur EDR & \textbf{FastAPI} (Python) — asynchrone, performant, bonne doc \\
    Agent EDR & \textbf{Python + PyInstaller} — binaire standalone, multi-OS \\
    Site client & \textbf{PHP 8+ / Apache} — simple, maîtrisé, suffisant \\
    Base clients & \textbf{MySQL} — fiable, bien supporté par PHP \\
    Base EDR & \textbf{PostgreSQL} — robuste, concurrence, fonctionnalités avancées \\
    Chatbot & \textbf{PHP → Claude API} — proxy léger, pas de framework \\
    Conteneurs & \textbf{Docker + GHCR} — multi-arch (amd64 + arm64) \\
    CI/CD & \textbf{GitHub Actions (homemade)} — pipeline custom complète \\
    Monitoring & \textbf{Prometheus / Grafana / Loki} — stack open-source standard \\
    \bottomrule
  \end{tabularx}

  \vspace{0.3cm}
  \begin{exampleblock}{Pourquoi pas Kubernetes ?}
    Trop lourd pour 4 conteneurs. Docker Compose suffit amplement.
    K3s reste une évolution possible si le projet passe à l'échelle.
  \end{exampleblock}
\end{frame}
\note{
  Chaque techno a été choisie pour sa pertinence par rapport à la taille du projet.
  FastAPI plutôt que Flask : meilleures perf, typage automatique avec Pydantic.
  PHP plutôt que React/Angular : pas besoin d'un framework lourd pour 5 pages.
  PostgreSQL plutôt que MySQL pour l'EDR : les événements nécessitent des requêtes avancées.
  CI/CD 100\% maison : pas de templates pré-construits, on a tout écrit.
}

% --- Slide 5 : Architecture Globale ---
\begin{frame}{Architecture Globale — Vue d'Ensemble}
  \begin{columns}[t]
    \column{0.48\textwidth}
    \textbf{4 composants}
    \begin{itemize}
      \bitem \textbf{Serveur EDR} : FastAPI, API REST, scan VT, tickets GLPI
      \bitem \textbf{Agent EDR} : Python, PyInstaller, FIM, Heartbeat
      \bitem \textbf{Site client} : PHP, MySQL, dashboard, GLPI SSO
      \bitem \textbf{Chatbot} : PHP, Claude API, rate limiting
    \end{itemize}
    \column{0.48\textwidth}
    \textbf{Flux de données}
    \begin{itemize}
      \bitem Agent $\rightarrow$ EDR API $\rightarrow$ VirusTotal / GLPI
      \bitem Client $\rightarrow$ Site PHP $\rightarrow$ MySQL / GLPI
      \bitem Chatbot $\rightarrow$ Claude API
      \bitem Git push $\rightarrow$ CI $\rightarrow$ GHCR $\rightarrow$ Dokploy $\rightarrow$ Prod
    \end{itemize}
  \end{columns}
  \vspace{0.3cm}
  \begin{alertblock}{3 dépôts Git (sous-modules)}
    \texttt{github.com/Lockbits-ESGI/main} (orchestrateur + CI/CD + docs) \\
    \texttt{edr} et \texttt{site\_lockbits} en sous-modules — mis à jour automatiquement toutes les heures
  \end{alertblock}
\end{frame}
\note{
  L'architecture est simple mais complète : 4 composants bien séparés, communiquant via API REST.
  Les sous-modules Git permettent de développer chaque composant indépendamment.
  Un workflow cron met à jour les sous-modules toutes les heures automatiquement.
}

% --- Slide 6 : CI/CD Pipeline ---
\begin{frame}{CI/CD — Pipeline 100\% Homemade}
  \begin{enumerate}
    \item \textbf{Push} sur \texttt{develop} ou \texttt{main}
    \item \textbf{Lint} : Ruff (Python), PHP CS Fixer
    \item \textbf{Tests} : pytest (EDR), PHPUnit (Site)
    \item \textbf{Build multi-arch} : Docker buildx (amd64 + arm64)
    \item \textbf{Push} vers GitHub Container Registry
    \item \textbf{Scan Trivy} : blocage si CRITICAL ou HIGH
    \item \textbf{Webhook} $\rightarrow$ Dokploy $\rightarrow$ déploiement auto
  \end{enumerate}
  \vspace{0.3cm}
  \textbf{Stratégie de déploiement}
  \begin{itemize}
    \bitem \texttt{main} $\rightarrow$ \texttt{:latest} $\rightarrow$ \textbf{production}
    \bitem \texttt{develop} $\rightarrow$ \texttt{:develop} $\rightarrow$ \textbf{pré-production}
  \end{itemize}
  \begin{exampleblock}{$\sim$5 min du push au déploiement — zéro intervention manuelle}
  \end{exampleblock}
\end{frame}
\note{
  Pipeline entièrement écrite à la main — pas un seul template GitHub Actions copié.
  Le build multi-arch prend le plus de temps (build pour amd64 ET arm64).
  Trivy scanne les vulnérabilités : si CRITICAL, le déploiement est bloqué.
  Dokploy reçoit un webhook et redéploie automatiquement le conteneur.
}

% --- Slide 7 : Hébergement & Infrastructure ---
\begin{frame}{Hébergement \& Infrastructure}
  \textbf{Serveur Proxmox personnel}
  \begin{itemize}
    \bitem Auto-hébergé chez Cléry, accessible 24/7
    \bitem \textbf{LXC Dokploy} : plateforme de déploiement automatisé
    \bitem \textbf{LXC cloudflared} : tunnel Cloudflare Zero Trust
    \bitem \textbf{Conteneurs Docker} : EDR, Site, Chatbot, Monitoring
  \end{itemize}
  \vspace{0.2cm}
  \textbf{Cloudflare Zero Trust}
  \begin{itemize}
    \bitem Connexion \textbf{sortante uniquement} — aucun port ouvert
    \bitem HTTPS automatique, protection DDoS, AuthZero
  \end{itemize}
  \vspace{0.2cm}
  \textbf{Services en ligne}
  \begin{itemize}
    \bitem \texttt{lockbits.pro} — Portail client
    \bitem \texttt{edr.lockbits.pro} — API EDR
    \bitem \texttt{chat.lockbits.pro} — Chatbot
  \end{itemize}
\end{frame}
\note{
  Tout tourne sur un vrai serveur Proxmox à la maison — pas de cloud.
  Cloudflare Zero Trust = le serveur n'expose aucun port, tout passe par le tunnel.
  C'est plus sûr car même si le serveur est compromis, l'attaquant ne peut pas ouvrir de ports entrants.
}

% --- Slide 8 : Monitoring & Observabilité ---
\begin{frame}{Monitoring \& Observabilité}
  \begin{columns}[t]
    \column{0.48\textwidth}
    \textbf{Stack PLA (Prometheus + Loki + Alloy)}
    \begin{itemize}
      \bitem \textbf{Prometheus} : métriques EDR (agents, événements, scans)
      \bitem \textbf{Loki} : logs centralisés (serveur, site, agent)
      \bitem \textbf{Alloy} : agent de collecte
      \bitem \textbf{Grafana} : dashboards + alerting
    \end{itemize}
    \column{0.48\textwidth}
    \textbf{Métriques clés}
    \begin{itemize}
      \bitem Agents connectés / actifs / silencieux
      \bitem Événements par minute
      \bitem Santé API, statut BDD
      \bitem RAM, CPU, disque
    \end{itemize}
  \end{columns}
  \vspace{0.3cm}
  \begin{alertblock}{Sécurité intégrée}
    Trivy scan sur chaque build — Secrets dans .env + GitHub Secrets — CORS + rate limiting
  \end{alertblock}
\end{frame}
\note{
  L'observabilité est essentielle : on doit savoir en temps réel si un agent est down ou si une menace est détectée.
  Prometheus scrape les métriques toutes les 15s.
  Grafana permet de visualiser l'état du système d'un coup d'œil.
}

% --- Slide 9 : Difficultés & Solutions (Cléry) ---
\begin{frame}{Difficultés Rencontrées — Infrastructure}
  \begin{description}
    \item[\textbf{Multi-architecture Docker}] \hfill \\
      \textbf{Problème} : Les images doivent tourner sur amd64 (serveur) ET arm64 (Mac M1/M2, Raspberry Pi). \\
      \textbf{Solution} : Utilisation de \texttt{docker buildx} avec \texttt{--platform linux/amd64,linux/arm64}. \\
      \textit{Résultat} : Une seule image GHCR qui fonctionne partout.
    \item[\textbf{Cloudflare Tunnel instable}] \hfill \\
      \textbf{Problème} : Le tunnel cloudflared tombait parfois, rendant les services inaccessibles. \\
      \textbf{Solution} : Mise en place d'un script de heartbeat + redémarrage automatique via systemd. \\
      \textit{Résultat} : Zéro downtime depuis la mise en place.
    \item[\textbf{Synchronisation sous-modules Git}] \hfill \\
      \textbf{Problème} : Les sous-modules n'étaient pas à jour sur le serveur de prod. \\
      \textbf{Solution} : Workflow GitHub Actions cron toutes les heures + script de pull auto sur le serveur. \\
      \textit{Résultat} : Déploiement toujours à jour.
  \end{description}
\end{frame}
\note{
  Difficultés techniques rencontrées en cours de route.
  Chaque problème a été résolu avec une solution pragmatique — pas de bidouillage.
  Le multi-arch était indispensable : le serveur Proxmox est en amd64, mais les devs sont sur Mac ARM.
}

% ========================================================================
% SECTION SWANN — EDR Server & Agent
% ========================================================================
\section{Swann — EDR Server \& Agent}

% --- Slide 10 : EDR Server Architecture ---
\begin{frame}{EDR Server — Architecture}
  \textbf{Un serveur FastAPI (Python) conçu pour la détection}
  \begin{itemize}
    \bitem API REST asynchrone (Uvicorn)
    \bitem Endpoints : \texttt{/health}, \texttt{/api/v1/stats}, \texttt{/api/v1/agents}, \texttt{/api/v1/events}
    \bitem Authentification : API key partagée avec l'agent
    \bitem Rate limiting + CORS configurable
  \end{itemize}
  \vspace{0.3cm}
  \textbf{Structure du code}
  \begin{itemize}
    \bitem \texttt{main.py} — Point d'entrée, configuration Uvicorn
    \bitem \texttt{api.py} — Routes et endpoints REST
    \bitem \texttt{schemas.py} — Modèles Pydantic (validation automatique)
    \bitem \texttt{metrics.py} — Métriques Prometheus
    \bitem \texttt{virus\_total.py} — Scan asynchrone VirusTotal
    \bitem \texttt{glpi.py} — Création automatique de tickets GLPI
  \end{itemize}
\end{frame}
\note{
  FastAPI = moderne, rapide, asynchrone.
  Les schémas Pydantic valident automatiquement les données entrantes/sortantes.
  L'API key est partagée entre serveur et agent, changée en production.
}

% --- Slide 11 : EDR Server — Fonctionnalités Clés ---
\begin{frame}{EDR Server — Détection \& Intégrations}
  \begin{block}{VirusTotal — Scan automatique}
    \begin{itemize}
      \bitem Fichier suspect soumis par l'agent $\rightarrow$ hash calculé $\rightarrow$ requête VT
      \bitem \textbf{Worker asynchrone} : pas de blocage du thread principal
      \bitem Résultat : score de réputation (0-100), moteurs de détection déclenchés
      \bitem Si score $>$ seuil $\rightarrow$ alerte immédiate + ticket GLPI
    \end{itemize}
  \end{block}
  \vspace{0.2cm}
  \begin{block}{GLPI — Ticketing automatique}
    \begin{itemize}
      \bitem Événement \textbf{CRITICAL} $\rightarrow$ création automatique d'un ticket GLPI
      \bitem Contenu : type de menace, agent, horodatage, hash du fichier
      \bitem Ticket visible dans le portail client en temps réel
      \bitem Authentification OAuth2 entre EDR et GLPI
    \end{itemize}
  \end{block}
  \vspace{0.2cm}
  \begin{exampleblock}{Parcours complet : Agent détecte $\rightarrow$ Serveur analyse $\rightarrow$ VT confirme $\rightarrow$ GLPI tickette $\rightarrow$ Client voit}
  \end{exampleblock}
\end{frame}
\note{
  VirusTotal agrège 70+ moteurs antivirus. Si le fichier est déjà connu, réponse en cache en < 1s.
  Le worker asynchrone évite de bloquer le serveur pendant le scan VT.
  GLPI est le back-office : l'utilisateur n'y accède pas, tout passe par le portail.
}

% --- Slide 12 : EDR Server — Base de Données ---
\begin{frame}{EDR Server — Persistance des Données}
  \textbf{PostgreSQL en production — SQLite en dev}
  \begin{itemize}
    \bitem PostgreSQL choisi pour : concurrence, transactions, requêtes avancées
    \bitem SQLite : déploiement local léger, pas de setup BDD
  \end{itemize}
  \vspace{0.3cm}
  \textbf{Tables principales}
  \begin{itemize}
    \bitem \texttt{agents} : hostname, OS, version, IP, last\_seen, statut (actif/silencieux/inactif)
    \bitem \texttt{events} : type (file/modified/process/connection), sévérité, timestamp, hostname
    \bitem \texttt{scans} : file\_hash, vt\_score, vt\_link, status (pending/done/error), agent
    \bitem \texttt{heartbeats} : agent, timestamp, status, uptime
  \end{itemize}
\end{frame}
\note{
  Choix motivé par le besoin : PostgreSQL pour la fiabilité en production, SQLite pour la simplicité en dev.
  Les schémas sont versionnés dans le dépôt Git.
  Les index sont optimisés pour les requêtes de recherche d'événements par agent et par période.
}

% --- Slide 13 : EDR Agent ---
\begin{frame}{EDR Agent — Architecture \& Fonctionnalités}
  \begin{columns}[t]
    \column{0.48\textwidth}
    \textbf{Agent Python $\rightarrow$ Binaire PyInstaller}
    \begin{itemize}
      \bitem Compilé standalone : Linux, macOS, Windows
      \bitem Taille : $\sim$15 Mo
      \bitem Aucune dépendance Python requise sur le poste client
    \end{itemize}
    \vspace{0.3cm}
    \textbf{Modules internes}
    \begin{itemize}
      \bitem \texttt{main.py} — Orchestrateur
      \bitem \texttt{heartbeat.py} — Signal de vie (30s)
      \bitem \texttt{sender.py} — Communication HTTP
      \bitem \texttt{fim.py} — File Integrity Monitoring
      \bitem \texttt{collector.py} — Infos système
      \bitem \texttt{hasher.py} — SHA-256
      \bitem \texttt{vt\_checker.py} — Scan VT
    \end{itemize}
    \column{0.48\textwidth}
    \textbf{2 modes de fonctionnement}
    \begin{itemize}
      \bitem \textbf{Scan} (one-shot) : analyse et rapport
      \bitem \textbf{Monitor} (daemon) : continu + heartbeat
    \end{itemize}
    \vspace{0.3cm}
    \textbf{Fonctionnalités de détection}
    \begin{itemize}
      \bitem Scan processus en cours
      \bitem Détection connexions réseau anormales
      \bitem \textbf{FIM} : SHA-256, fichiers critiques (/etc/passwd, binaires)
      \bitem Vérification réputation VirusTotal
    \end{itemize}
    \vspace{0.3cm}
    \textbf{Déploiement one-liner}
    \begin{itemize}
      \bitem \texttt{curl -sSL ... | sudo bash}
      \bitem Auto-détection OS + architecture
      \bitem Installation en $\sim$10 secondes
    \end{itemize}
  \end{columns}
\end{frame}
\note{
  L'agent est le composant le plus critique : il tourne sur le poste du client.
  PyInstaller permet de livrer un binaire sans exposer le code source.
  Le FIM (File Integrity Monitoring) détecte les modifications de fichiers systèmes — essentiel pour la sécurité.
  Installation one-liner : une commande, pas de questions, ça marche.
}

% --- Slide 14 : Communication Agent-Serveur ---
\begin{frame}{Communication Agent $\leftrightarrow$ Serveur}
  \textbf{Protocole}
  \begin{itemize}
    \bitem HTTP POST, payload JSON
    \bitem Authentification : header \texttt{X-API-Key}
    \bitem Payload : type d'événement, données, horodatage, hash
  \end{itemize}
  \vspace{0.3cm}
  \textbf{Cycle de vie d'un événement}
  \begin{enumerate}
    \item Agent détecte un fichier suspect $\rightarrow$ calcule SHA-256
    \item Envoie au serveur via API REST
    \item Serveur vérifie le hash (déjà vu ? connu VT ?)
    \item Si nouveau : scan VirusTotal asynchrone
    \item Résultat : bénin / suspect (alerte) / critique (ticket GLPI)
    \item Agent reçoit la réponse : OK / ALERT / QUARANTINE
  \end{enumerate}
\end{frame}
\note{
  Protocole simple mais efficace : HTTP + JSON + API key.
  L'agent attend une instruction en retour : agir ou non.
  En mode monitor, les événements sont envoyés en continu.
}

% --- Slide 15 : Difficultés & Solutions (Swann) ---
\begin{frame}{Difficultés Rencontrées — EDR}
  \begin{description}
    \item[\textbf{Compilation PyInstaller multi-OS}] \hfill \\
      \textbf{Problème} : PyInstaller ne produit un binaire que pour l'OS hôte. Impossible de compiler Windows depuis Linux. \\
      \textbf{Solution} : Build séparé par OS via GitHub Actions matrix, chaque job compile sur l'OS cible. \\
      \textit{Résultat} : 3 binaires (Linux, macOS, Windows) générés automatiquement à chaque release.
    \item[\textbf{Déploiement de l'agent chez le client}] \hfill \\
      \textbf{Problème} : L'agent doit s'installer sans interaction, sur n'importe quel OS. \\
      \textbf{Solution} : Script \texttt{install-agent.sh} qui détecte l'OS, télécharge le bon binaire et configure systemd/launchd. \\
      \textit{Résultat} : \texttt{curl ... | sudo bash} — 10 secondes, zéro question.
    \item[\textbf{Intégration VirusTotal limitée}] \hfill \\
      \textbf{Problème} : API VirusTotal gratuite limitée à 4 requêtes/minute. \\
      \textbf{Solution} : File d'attente locale + cache des hash déjà scannés (évite les doublons). \\
      \textit{Résultat} : Optimisation des requêtes, pas de dépassement de quota.
  \end{description}
\end{frame}
\note{
  Difficultés techniques concrètes sur l'EDR.
  PyInstaller : on a perdu du temps à chercher une solution cross-platform, la vraie solution c'est GitHub Actions matrix.
  Le one-liner : beaucoup d'itérations pour gérer tous les cas (curl absent, OS non supporté, permissions).
  VirusTotal gratuit : contrainte réelle, contournée par le cache et la file d'attente.
}

% ========================================================================
% SECTION SOLTANE — Site Client
% ========================================================================
\section{Soltane — Site Client}

% --- Slide 16 : Site Client Architecture ---
\begin{frame}{Site Client — Architecture}
  \textbf{Portail web complet pour les clients LockBits}
  \begin{itemize}
    \bitem \textbf{PHP 8+ / Apache} — stack classique, fiable, bien maîtrisée
    \bitem \textbf{MySQL} — persistance des données clients et tickets
    \bitem \textbf{Tailwind CSS} — design moderne sans framework JS lourd
    \bitem \textbf{Docker} — image GHCR, multi-arch, déploiement automatisé
  \end{itemize}
  \vspace{0.3cm}
  \textbf{Fonctionnalités du portail}
  \begin{itemize}
    \bitem \textbf{Authentification} — Inscription, connexion, profil, rôles (client/admin)
    \bitem \textbf{Dashboard} — Vue d'ensemble : alertes, statut abonnement, tickets récents
    \bitem \textbf{Tickets} — Consultation et suivi des tickets GLPI synchronisés
    \bitem \textbf{Téléchargement} — Accès au binaire de l'agent EDR
    \bitem \textbf{Support} — Chatbot intégré directement dans le portail
  \end{itemize}
\end{frame}
\note{
  Le site est l'interface principale du client.
  PHP choisi pour sa simplicité et la maîtrise de l'équipe.
  Tailwind CSS permet un design propre sans écrire de CSS.
  Pas de JavaScript lourd : le site reste rapide et accessible.
}

% --- Slide 17 : Site Client — Authentification, BDD, SSO ---
\begin{frame}{Site Client — Sécurité \& Intégrations}
  \begin{columns}[t]
    \column{0.48\textwidth}
    \textbf{Authentification sécurisée}
    \begin{itemize}
      \bitem Mots de passe hachés en \textbf{bcrypt}
      \bitem Sessions PHP sécurisées, protection CSRF
      \bitem Requêtes préparées (PDO) — pas d'injection SQL
      \bitem Validation des entrées côté serveur
    \end{itemize}
    \vspace{0.3cm}
    \textbf{Base MySQL}
    \begin{itemize}
      \bitem \texttt{clients} : email, password (hash), rôle, abonnement
      \bitem \texttt{tickets} : ref GLPI, titre, statut, priorité
      \bitem \texttt{metrics} : connexions, téléchargements
    \end{itemize}
    \column{0.48\textwidth}
    \textbf{GLPI SSO — Connexion unique}
    \begin{itemize}
      \bitem \textbf{OAuth2} : connexion via compte GLPI
      \bitem Une seule session : GLPI + Portail
      \bitem Tickets synchronisés en temps réel
    \end{itemize}
    \vspace{0.3cm}
    \textbf{Parcours utilisateur}
    \begin{enumerate}
      \item Client se connecte sur \texttt{lockbits.pro}
      \item Voit son dashboard + ses tickets GLPI
      \item Une alerte EDR critique arrive
      \item Ticket créé automatiquement dans GLPI
      \item Visible dans le portail en \textbf{temps réel}
    \end{enumerate}
  \end{columns}
\end{frame}
\note{
  La sécurité est une priorité : bcrypt, PDO, CSRF, validation.
  Le SSO GLPI évite la double authentification.
  Le parcours client est fluide : pas de déconnexion entre GLPI et le portail.
  Les tickets sont synchronisés via API REST GLPI.
}

% --- Slide 18 : Site Client — Déploiement ---
\begin{frame}{Site Client — Déploiement \& Métriques}
  \textbf{Déploiement automatisé}
  \begin{itemize}
    \bitem Pipeline CI/CD identique aux autres composants
    \bitem Dockerfile Apache + PHP, image GHCR
    \bitem Déploiement via Dokploy : preprod puis prod
  \end{itemize}
  \vspace{0.2cm}
  \textbf{Sauvegardes}
  \begin{itemize}
    \bitem \textbf{Backup MySQL automatisé} : script \texttt{backup.sh}
    \bitem Rotation sur 7 jours, téléchargement sécurisé
    \bitem Testé et vérifié (restauration possible)
  \end{itemize}
  \vspace{0.2cm}
  \textbf{Métriques collectées}
  \begin{itemize}
    \bitem Connexions, inscriptions, téléchargements agent
    \bitem Tickets créés/résolus/en attente
    \bitem Temps de réponse, erreurs serveur
    \bitem Pages vues, actions utilisateur
  \end{itemize}
\end{frame}
\note{
  Même pipeline CI/CD que l'EDR — chaque push testé et déployé.
  Les backups MySQL sont automatisés : pas de risque de perte de données.
  Les métriques permettent de suivre l'adoption du service.
}

% --- Slide 19 : Difficultés & Solutions (Soltane) ---
\begin{frame}{Difficultés Rencontrées — Site Client}
  \begin{description}
    \item[\textbf{Synchronisation GLPI temps réel}] \hfill \\
      \textbf{Problème} : Les tickets GLPI devaient être visibles immédiatement dans le portail. \\
      \textbf{Solution} : Webhook GLPI $\rightarrow$ mise à jour auto côté portail + polling toutes les 30s. \\
      \textit{Résultat} : Synchronisation quasi-temps réel (délai < 30s).
    \item[\textbf{Dockerisation PHP/Apache}] \hfill \\
      \textbf{Problème} : Apache nécessite des modules spécifiques (mod\_rewrite, mod\_ssl) non présents dans l'image de base PHP. \\
      \textbf{Solution} : Dockerfile personnalisé avec installation des modules Apache + configuration PHP optimisée. \\
      \textit{Résultat} : Image légère ($\sim$200 Mo) avec tout le nécessaire.
    \item[\textbf{Protection CSRF sur les formulaires}] \hfill \\
      \textbf{Problème} : Les formulaires de connexion et tickets étaient vulnérables aux attaques CSRF. \\
      \textbf{Solution} : Implémentation de tokens CSRF avec génération aléatoire + validation côté serveur. \\
      \textit{Résultat} : Protection complète contre les attaques cross-site.
  \end{description}
\end{frame}
\note{
  Difficultés rencontrées sur le site.
  La synchro GLPI a été le plus complexe : il fallait que le client voie ses tickets sans rafraîchir la page.
  La Dockerisation PHP/Apache : les images officielles ne sont pas toujours adaptées, il faut les personnaliser.
  CSRF : attaque classique, protection standard mais indispensable.
}

% ========================================================================
% SECTION QUENTIN — Chatbot \& Retour d'Expérience
% ========================================================================
\section{Quentin — Chatbot \& Retour d'Expérience}

% --- Slide 20 : Chatbot Architecture ---
\begin{frame}{Chatbot — Architecture}
  \textbf{Assistant virtuel basé sur Claude API (Anthropic)}
  \begin{itemize}
    \bitem \textbf{Proxy PHP} : \texttt{chat.php}, point d'entrée unique
    \bitem Intégré au site client (fenêtre de chat)
    \bitem Aucun framework — fichier unique, léger
  \end{itemize}
  \vspace{0.3cm}
  \textbf{Fonctionnement}
  \begin{enumerate}
    \item Utilisateur envoie un message
    \item PHP valide et nettoie l'entrée (sécurité)
    \item Requête forwardée à Claude API avec system prompt verrouillé
    \item Claude génère une réponse $\rightarrow$ PHP la vérifie $\rightarrow$ l'affiche
  \end{enumerate}
  \vspace{0.3cm}
  \textbf{Fonctionnalités}
  \begin{itemize}
    \bitem Questions sur LockBits : services, abonnement, utilisation
    \bitem Aide au diagnostic : interprétation des alertes EDR
    \bitem Historique de conversation (session)
    \bitem Réponses naturelles grâce à l'IA Claude
  \end{itemize}
\end{frame}
\note{
  Le chatbot permet aux clients d'obtenir des réponses sans solliciter l'équipe.
  Claude API génère des réponses en langage naturel.
  Le system prompt limite le chatbot au contexte LockBits — il ne répond pas à des questions dangereuses.
}

% --- Slide 21 : Chatbot — Sécurité & Difficultés ---
\begin{frame}{Chatbot — Sécurité \& Difficultés}
  \begin{columns}[t]
    \column{0.48\textwidth}
    \textbf{Sécurité}
    \begin{itemize}
      \bitem Rate limiting : \textbf{100 requêtes/jour/IP}, fichier compteur JSON
      \bitem Clé API Claude dans \texttt{.env}, jamais exposée
      \bitem Validation des entrées : pas d'injection, pas de code
      \bitem System prompt verrouillé : périmètre LockBits uniquement
      \bitem Messages anonymisés avant envoi à Claude
    \end{itemize}
    \column{0.48\textwidth}
    \textbf{Difficultés rencontrées}
    \begin{itemize}
      \bitem \textbf{Clé API exposée} : au début, la clé était dans le code JS.\\
            \textit{Solution} : Proxy PHP côté serveur, plus jamais exposée.
      \bitem \textbf{Rate limiting agressif} : Claude coupe après 100 requêtes/min.\\
            \textit{Solution} : Compteur quotidien, pas par minute.
      \bitem \textbf{Coût API} : chaque requête coûte $\sim$0,01\$.\\
            \textit{Solution} : Limite à 100/jour, coût maîtrisé ($<$ 10\$/mois).
    \end{itemize}
  \end{columns}
\end{frame}
\note{
  La sécurité du chatbot est cruciale : il ne doit pas divulguer d'informations sensibles.
  Le proxy PHP évite d'exposer la clé API côté client.
  Le rate limiting évite les abus et maîtrise les coûts.
  Les messages sont anonymisés : aucune donnée personnelle transmise à Claude.
}

% --- Slide 22 : Résultats ---
\begin{frame}{Résultats \& Services en Ligne}
  \begin{exampleblock}{Services accessibles publiquement}
    \centering
    \begin{tabularx}{\textwidth}{@{}l X@{}}
      \toprule
      \textbf{Service} & \textbf{URL} \\
      \midrule
      Portail client & \texttt{https://lockbits.pro} \\
      API EDR & \texttt{https://edr.lockbits.pro} \\
      Pré-production & \texttt{https://preprod-edr.lockbits.pro} \\
      Chatbot & \texttt{https://chat.lockbits.pro} \\
      \bottomrule
    \end{tabularx}
  \end{exampleblock}
  \vspace{0.3cm}
  \textbf{Chiffres clés}
  \begin{itemize}
    \bitem Pipeline CI \textbf{verte en permanence} — chaque PR testée avant merge
    \bitem Déploiement : \textasciitilde5 min du push à la production
    \bitem Agent EDR : installation one-liner, 3 OS supportés
    \bitem Code 100\% open-source : \texttt{github.com/Lockbits-ESGI}
    \bitem Dossier d'architecture LaTeX complet (\textasciitilde2800 lignes)
  \end{itemize}
\end{frame}
\note{
  Résultats concrets du projet : tout est en ligne, accessible, fonctionnel.
  La pipeline CI est verte — preuve que le projet est industrialisé.
  Le dossier d'architecture LaTeX est livré avec le code source.
}

% --- Slide 23 : Retour d'Expérience ---
\section{Retour d'Expérience — Toute l'Équipe}
\begin{frame}{Retour d'Expérience — Ce que le Projet Nous a Apporté}
  \begin{block}{Apports du projet annuel}
    \begin{itemize}
      \bitem \textbf{Autonomie} : gérer un projet de A à Z sans encadrement technique quotidien
      \bitem \textbf{Travail d'équipe} : coordination entre 4 personnes avec des compétences complémentaires
      \bitem \textbf{Deadline réelle} : livrer pour juillet, gérer son temps et ses priorités
      \bitem \textbf{De l'idée à la prod} : coder, builder, déployer, maintenir — tout le cycle
    \end{itemize}
  \end{block}
  \vspace{0.2cm}
  \begin{block}{Compétences acquises par domaine}
    \begin{itemize}
      \bitem \textbf{Cléry} : CI/CD, Docker multi-arch, Proxmox, Cloudflare Zero Trust, orchestration Git
      \bitem \textbf{Swann} : FastAPI, PyInstaller, VirusTotal API, architecture agent-serveur, monitoring
      \bitem \textbf{Soltane} : PHP 8+, MySQL, GLPI API, OAuth2 SSO, Tailwind CSS, Dockerisation Apache
      \bitem \textbf{Quentin} : Claude API, sécurité applicative, rate limiting, proxy PHP, intelligence artificielle
    \end{itemize}
  \end{block}
\end{frame}
\note{
  Le projet annuel nous a apporté une expérience concrète de A à Z.
  Chacun a développé des compétences spécifiques à son domaine.
  L'autonomie a été le plus grand défi : pas de professeur pour guider, il fallait chercher les solutions soi-même.
}

% --- Slide 24 : Pistes d'Amélioration ---
\begin{frame}{Pistes d'Amélioration}
  \textbf{Techniques}
  \begin{itemize}
    \bitem \textbf{Couverture de tests} : viser 80\%+ sur l'EDR (actuellement $\sim$50\%)
    \bitem \textbf{Dashboard Grafana pro} : dashboards par composant avec alerting poussé
    \bitem \textbf{Kubernetes (K3s)} : si le projet passe à l'échelle, orchestration complète
  \end{itemize}
  \vspace{0.2cm}
  \textbf{Fonctionnalités}
  \begin{itemize}
    \bitem \textbf{OAuth2 complet} : Google, Microsoft, GitHub — en plus de GLPI
    \bitem \textbf{SIEM} : Wazuh ou ELK pour la corrélation d'événements avancée
    \bitem \textbf{Interface admin web EDR} : gestion des agents, alertes, rapports
  \end{itemize}
  \vspace{0.2cm}
  \textbf{Organisationnelles}
  \begin{itemize}
    \bitem \textbf{Plus de tests utilisateurs} : recueillir des retours avant la mise en prod
    \bitem \textbf{Documentation utilisateur} : guide d'utilisation du portail et de l'agent
    \bitem \textbf{Méthodologie plus formelle} : sprints mieux cadrés, daily plus réguliers
  \end{itemize}
\end{frame}
\note{
  On a livré un projet fonctionnel, mais il y a toujours des pistes d'amélioration.
  Les tests : c'est ce qu'on améliorerait en priorité avec plus de temps.
  K3s : pas nécessaire aujourd'hui mais logique si le projet grandit.
  Documentation utilisateur : on a la doc technique (LaTeX) mais pas de guide utilisateur simple.
}

% --- Slide 25 : Merci ---
\begin{frame}{Merci}
  \vspace{1.5em}
  \centering
  {\Huge\color{primaryblue} Des questions ?}
  \vspace{1em}
  \rule{0.5\textwidth}{0.5pt}
  \vspace{1.5em}
  {\Large\url{https://github.com/Lockbits-ESGI}}
  \vspace{0.5em}
  \begin{itemize}
    \bitem Dossier d'architecture : \texttt{docs/latex/} (LaTeX, \textasciitilde2800 lignes)
    \bitem Présentation : \texttt{docs/latex/presentation.tex}
    \bitem Code source : \texttt{github.com/Lockbits-ESGI/main}
  \end{itemize}
  \vspace{1em}
  \small{\textit{LockBits — EDR Open-Source \& Portail Client — ESGI 2025-2026}}
\end{frame}
\note{
  Merci au jury. Proposer une démo en direct si le temps le permet.
  Tous les services sont en ligne, le code est open-source — on peut tout montrer.
  Chacun répond sur son domaine : Cléry (infra/CI/CD), Swann (EDR), Soltane (site), Quentin (chatbot/retour).
}

\end{document}
